% The Explanation Impossibility: A Universal Trilemma for Interpreting
% Underspecified Systems
%
% NeurIPS 2026 submission — pages 1–5 (first half)

\documentclass{article}
\usepackage[preprint]{neurips_2026}
% natbib loaded by neurips_2026.sty
\usepackage{amsmath,amssymb,amsthm}
\usepackage{booktabs}
\usepackage{enumitem}

\DeclareMathOperator{\Var}{Var}
\newtheorem{theorem}{Theorem}
\newtheorem{corollary}[theorem]{Corollary}
\newtheorem{proposition}[theorem]{Proposition}
\newtheorem{definition}[theorem]{Definition}
\newtheorem{remark}[theorem]{Remark}

\title{The Explanation Impossibility:\\A Universal Trilemma for Interpreting Underspecified Systems}

\author{
  Anonymous Authors
}

\begin{document}

\maketitle

%% ===================================================================
%% ABSTRACT
%% ===================================================================
\begin{abstract}
We prove that no explanation of an underspecified system can be simultaneously
\emph{faithful} (consistent with the system's native explanation),
\emph{stable} (invariant across equivalent configurations), and
\emph{decisive} (preserving genuine disagreements).
For binary explanation problems---including mechanistic interpretability and
SHAP sign attribution---the impossibility strengthens to a \emph{bilemma}:
faithfulness and stability alone are jointly unachievable.
The result is universal: a single four-line proof, with zero domain-specific
axioms, covers every instance we tested.
We classify which problems admit resolutions and which do not:
the algebraic structure of the explanation space (presence or absence of a
neutral element) determines the ceiling.
We instantiate the framework across mechanistic interpretability (circuit
identification), feature attribution (SHAP rankings), causal discovery
(edge orientation), and attention-based token importance, showing each is a
special case of the same theorem.
When a resolution exists, we characterize the optimal one via the
Hunt--Stein theorem: the Pitman estimator (group-orbit average).
The entire framework---357 theorems from 6 axioms, 0 \texttt{sorry}---is formalized in
Lean~4, with all binary explanation types (\texttt{SHAPSign},
\texttt{FeatureStatus}, \texttt{CounterfactualDir}) defined as
constructive inductive types and the mechanistic interpretability
Rashomon property \emph{derived} by computation (zero axioms).
\end{abstract}

%% ===================================================================
%% 1. INTRODUCTION
%% ===================================================================
\section{Introduction}
\label{sec:intro}

The search for circuits in neural networks is a cornerstone of AI safety.
Mechanistic interpretability aims to identify which internal components---neurons,
attention heads, subnetworks---implement a given behavior, with the ultimate goal
of verifying that a model's reasoning is aligned with human intent.
We prove that this search faces a fundamental mathematical constraint:
no circuit explanation can be simultaneously faithful to the network's internal
structure, stable across functionally equivalent networks, and decisive about
which decomposition is correct.
For binary circuit questions---``is this component part of the circuit or
not?''---the situation is strictly worse: faithfulness and stability alone
are jointly impossible, regardless of whether decisiveness is required.

This result is not specific to mechanistic interpretability.
The same impossibility governs every explanation method we tested:
SHAP feature rankings, causal edge orientations, attention-based token
importance, feature selection, and counterfactual directions.
All are instances of a single abstract theorem that requires
\emph{zero domain-specific axioms}---only the observation that underspecified
systems admit equivalent configurations with incompatible native explanations
(the Rashomon property).

The impossibility admits resolution for some problems but not others.
The key structural diagnostic is whether the explanation space contains a
\emph{neutral element}---a value compatible with every native explanation.
Feature rankings admit ties (a neutral element), so ensemble averaging
(DASH) resolves the trilemma.
Circuit decompositions do not admit a neutral element, so the bilemma is
inescapable; the only resolution is to enrich the explanation space by
reporting equivalence classes of circuits.
This classification is complete: the presence or absence of neutral and
committal elements determines exactly which pairs of axioms can be jointly
satisfied.

\paragraph{Contributions.}
\begin{enumerate}[leftmargin=*,label=(\arabic*)]
\item \textbf{Abstract framework and universal impossibility.}
  We define an \texttt{ExplanationSystem} $(\Theta, H, Y, \mathrm{observe},
  \mathrm{explain}, \mathrm{incompatible})$ and prove that no explanation
  $E : \Theta \to H$ of a system with the Rashomon property can be
  faithful, stable, and decisive (Theorem~\ref{thm:trilemma}).
  The proof is four lines and uses zero domain-specific axioms.

\item \textbf{The bilemma.}
  For maximally incompatible explanation spaces (compatible $\Rightarrow$
  equal), faithful + stable alone is impossible
  (Theorem~\ref{thm:bilemma})---strictly stronger than the trilemma.
  This applies to all binary explanation problems: SHAP sign, circuit
  decomposition, feature selection, and counterfactual direction
  (increase vs.\ decrease a feature).
  Each is formalized as a named bilemma instance
  (\texttt{shap\_sign\_bilemma}, \texttt{feature\_selection\_bilemma},
  \texttt{counterfactual\_bilemma}).

\item \textbf{Tightness classification.}
  We give a complete $2 \times 2$ classification
  (Table~\ref{tab:tightness}) showing which pairs of axioms are jointly
  achievable as a function of the explanation space's algebraic structure
  (neutral elements $\times$ committal elements).

\item \textbf{Instances.}
  We instantiate the framework across mechanistic interpretability
  (\S\ref{sec:mech-interp}), feature attribution (\S\ref{sec:attribution}),
  causal discovery (\S\ref{sec:causal}), and attention maps
  (\S\ref{sec:attention}), with additional instances in the supplement.

\item \textbf{Optimal resolution via Hunt--Stein.}
  When a resolution exists (neutral element present), the optimal
  explanation is the Pitman estimator---the average over the symmetry group
  orbit. For feature rankings, this is DASH (ensemble SHAP averaging).
  For causal discovery, this is the CPDAG.

\item \textbf{Lean~4 formalization.}
  The full framework is formalized across 58 files containing 357
  theorems and lemmas with 0~\texttt{sorry} and only 6 axioms
  (all infrastructure: type declarations and measure setup).
  The mechanistic interpretability Rashomon property is \emph{derived}
  from constructive finite types, not axiomatized.
  Three domain-specific inductive types---\texttt{SHAPSign},
  \texttt{FeatureStatus}, \texttt{CounterfactualDir}---are
  defined constructively with zero axioms.
\end{enumerate}

\paragraph{Scope separation.}
A companion paper~\citep{anonymous2026attribution} develops the
quantitative theory for feature rankings under collinearity: the
$1/(1{-}\rho^2)$ divergence rate, Spearman instability bounds, and
GBDT/Lasso/neural-net instantiations.
The present paper abstracts the core mechanism into a domain-independent
framework, proves the bilemma (absent from the companion),
and extends the impossibility to mechanistic interpretability,
causal discovery, attention maps, and counterfactual directions.

%% ===================================================================
%% 2. THE FRAMEWORK
%% ===================================================================
\section{The Framework}
\label{sec:framework}

\begin{definition}[Explanation System]
\label{def:explanation-system}
An \emph{explanation system} is a tuple $S = (\Theta, H, Y,
\mathrm{observe}, \mathrm{explain}, \mathrm{incompatible})$ where:
\begin{itemize}[leftmargin=*]
\item $\Theta$ is the \emph{configuration space} (e.g., trained models,
  network weight vectors, DAGs);
\item $H$ is the \emph{explanation space} (e.g., feature rankings, circuit
  decompositions, edge orientations);
\item $Y$ is the \emph{observable space} (e.g., predictions, I/O functions,
  conditional independence structures);
\item $\mathrm{observe} : \Theta \to Y$ maps each configuration to its
  observable behavior;
\item $\mathrm{explain} : \Theta \to H$ maps each configuration to its
  \emph{native} explanation (the ground-truth interpretation that the
  configuration's internals warrant);
\item $\mathrm{incompatible} : H \times H \to \mathrm{Prop}$ is an
  irreflexive relation on explanations capturing genuine disagreement.
\end{itemize}
\end{definition}

An \emph{explanation method} is any function $E : \Theta \to H$.
We require three properties:

\paragraph{Faithfulness.}
$E$ is \emph{faithful} if it never contradicts the system's native
explanation:
\[
  \forall\, \theta \in \Theta,\quad
  \neg\,\mathrm{incompatible}(E(\theta),\; \mathrm{explain}(\theta)).
\]
This is deliberately weak: $E(\theta)$ need not \emph{equal}
$\mathrm{explain}(\theta)$, only be compatible with it.
A faithful explanation may be vague (e.g., report a tie) but must not
actively contradict the model.

\paragraph{Stability.}
$E$ is \emph{stable} if equivalent configurations receive the same
explanation:
\[
  \forall\, \theta_1, \theta_2 \in \Theta,\quad
  \mathrm{observe}(\theta_1) = \mathrm{observe}(\theta_2)
  \;\Longrightarrow\;
  E(\theta_1) = E(\theta_2).
\]
This is the explanation analogue of the statistical principle that
functionally equivalent models should yield the same interpretation.
Without stability, explanations depend on arbitrary implementation details
(random seed, initialization, training order) rather than the system's
behavior.

\paragraph{Decisiveness.}
$E$ is \emph{decisive} if it preserves the incompatibility of native
explanations:
\[
  \forall\, \theta_1, \theta_2 \in \Theta,\quad
  \mathrm{incompatible}(\mathrm{explain}(\theta_1),\;
  \mathrm{explain}(\theta_2))
  \;\Longrightarrow\;
  \mathrm{incompatible}(E(\theta_1),\; E(\theta_2)).
\]
A decisive method does not smooth over genuine disagreements between
configurations.
When two models truly disagree about which feature is more important
(or which circuit implements a behavior), a decisive method reports
that disagreement rather than collapsing it.

\paragraph{The Rashomon property.}
An explanation system has the \emph{Rashomon property} if there exist
configurations $\theta_1, \theta_2 \in \Theta$ such that
$\mathrm{observe}(\theta_1) = \mathrm{observe}(\theta_2)$ and
$\mathrm{incompatible}(\mathrm{explain}(\theta_1),\;
\mathrm{explain}(\theta_2))$.
That is, the system admits equivalent configurations whose native
explanations genuinely disagree.
This is the mechanism behind model multiplicity: the same predictive
performance can arise from models with contradictory internal structures.

%% ===================================================================
%% 3. THE IMPOSSIBILITY
%% ===================================================================
\section{The Impossibility}
\label{sec:impossibility}

\begin{theorem}[The Explanation Trilemma]
\label{thm:trilemma}
Let $S = (\Theta, H, Y, \mathrm{observe}, \mathrm{explain},
\mathrm{incompatible})$ be an explanation system where
$\mathrm{incompatible}$ is irreflexive.
If $S$ has the Rashomon property, then no explanation method
$E : \Theta \to H$ can be simultaneously faithful, stable, and decisive.
\end{theorem}

\begin{proof}
Let $\theta_1, \theta_2$ be Rashomon witnesses:
$\mathrm{observe}(\theta_1) = \mathrm{observe}(\theta_2)$ and
$\mathrm{incompatible}(\mathrm{explain}(\theta_1),
\mathrm{explain}(\theta_2))$.
(1)~Decisive: $\mathrm{incompatible}(E(\theta_1), E(\theta_2))$.
(2)~Stable: $E(\theta_1) = E(\theta_2)$.
(3)~Substituting: $\mathrm{incompatible}(E(\theta_1), E(\theta_1))$.
(4)~This contradicts irreflexivity.
\end{proof}

\begin{theorem}[Tightness]
\label{thm:tightness}
The trilemma is tight: each pair of axioms is jointly satisfiable.
\end{theorem}
\begin{proof}[Proof sketch]
(F+S, drop D): Let $E$ be a constant function mapping every $\theta$ to a
  neutral element $h_0$ compatible with all native explanations. Then $E$ is
  trivially faithful and stable, but not decisive.
(F+D, drop S): Let $E = \mathrm{explain}$. Then $E$ is faithful and
  decisive by definition, but not stable (Rashomon witnesses get different
  explanations).
(S+D, drop F): Let $E$ be a constant function mapping every $\theta$ to
  some fixed $h^*$ that is incompatible with at least one native
  explanation. Then $E$ is stable and (vacuously) decisive on pairs that
  agree, but unfaithful.
\end{proof}

\begin{theorem}[The Bilemma]
\label{thm:bilemma}
Let $S$ be an explanation system with the Rashomon property, where $H$ is
\emph{maximally incompatible}: $\neg\,\mathrm{incompatible}(h_1, h_2)
\Rightarrow h_1 = h_2$.
Then no explanation method $E : \Theta \to H$ can be simultaneously
faithful and stable.
\end{theorem}
\begin{proof}
Since incompatibility is the complement of equality, faithfulness
($\neg\,\mathrm{incompatible}(E(\theta), \mathrm{explain}(\theta))$)
forces $E(\theta) = \mathrm{explain}(\theta)$ for all $\theta$.
But then $E = \mathrm{explain}$ is decisive (it inherits the
incompatibility structure of $\mathrm{explain}$).
So $E$ is faithful, stable, and decisive, contradicting
Theorem~\ref{thm:trilemma}.
\end{proof}

\paragraph{The tightness classification.}
Whether faithful + stable is achievable depends on the algebraic structure
of the explanation space $H$.
A \emph{neutral element} $h_0 \in H$ satisfies
$\neg\,\mathrm{incompatible}(h_0, h)$ for all $h \in H$ (it is compatible
with every native explanation).
A \emph{committal element} $h^* \in H$ satisfies
$\mathrm{incompatible}(h^*, h)$ for some $h \in H$ (it actively disagrees
with at least one native explanation).

\begin{table}[t]
\centering
\caption{Tightness classification. Which pairs of axioms are jointly
achievable depends on whether $H$ contains neutral and committal
elements. F+S+D is never achievable (the trilemma).}
\label{tab:tightness}
\begin{tabular}{@{}lcc@{}}
\toprule
& Neutral exists & No neutral \\
\midrule
\textbf{Committal exists} & F+D,\; F+S,\; S+D & F+D,\; S+D \\
\textbf{No committal}     & F+D,\; F+S         & F+D only \\
\bottomrule
\end{tabular}
\end{table}

The key insight: \emph{the explanation space's algebraic structure
determines the ceiling.}
Feature rankings have ties (neutral elements), so DASH achieves F+S.
Circuit decompositions have no neutral element (every decomposition commits
to a structure), so F+S is impossible---the bilemma is inescapable.
This yields a \emph{direction theorem}: coverage conflict (the absence of neutral elements) is
anti-monotone in $|H|$---enriching $H$ with neutral elements weakens the
impossibility, the opposite of Arrow's theorem where enlarging the alternative set
strengthens the impossibility by enabling preference cycles.
The diagnostic is actionable: before building an explanation method,
check whether $H$ has a neutral element.
If not, no method can be both faithful and stable, and the only resolution
is to enrich $H$ (e.g., report equivalence classes instead of individual
explanations).

\paragraph{Coverage conflict diagnostic.}
We formalize the diagnostic as a predicate \texttt{hasCoverageConflict}
(\texttt{BeyondBinary.lean}):
an explanation system has \emph{coverage conflict} if every element
$h \in H$ is incompatible with some native explanation value
($\forall\, h,\; \exists\, \theta,\;
\mathrm{incompatible}(h, \mathrm{explain}(\theta))$).
Coverage conflict implies no neutral element exists
(Lean-verified: \texttt{coverageConflict\_implies\_no\_neutral}),
so the bilemma applies.
Conversely, a neutral element destroys coverage conflict
(\texttt{neutral\_destroys\_coverageConflict}).
The tightness dichotomy---F+S achievable $\Leftrightarrow$ absence
of coverage conflict---is proved as a single theorem
(\texttt{tightness\_dichotomy}).

%% ===================================================================
%% 4. INSTANCES
%% ===================================================================
\section{Instances}
\label{sec:instances}

We instantiate the framework across four explanation domains.
Each requires only specifying the tuple $(\Theta, H, Y, \mathrm{observe},
\mathrm{explain}, \mathrm{incompatible})$ and verifying the Rashomon
property; the impossibility then follows from Theorem~\ref{thm:trilemma}
or~\ref{thm:bilemma} with no additional proof.

%% ----- Instance 1: Mechanistic Interpretability -----
\subsection{Instance 1: Mechanistic Interpretability}
\label{sec:mech-interp}

\paragraph{Setup.}
$\Theta$ is the space of neural network configurations (weight vectors);
$H$ is the space of circuit decompositions (which subnetwork implements a
given behavior);
$Y$ is the space of input-output functions.
Two networks are equivalent ($\mathrm{observe}(\theta_1) =
\mathrm{observe}(\theta_2)$) if they compute the same function.
Two decompositions are incompatible if they identify different components
as implementing the behavior.

\paragraph{The bilemma applies.}
Circuit decomposition is a binary question at the component level:
each neuron or attention head is either part of the circuit or not.
The explanation space is maximally incompatible
($\neg\,\mathrm{incompatible}(d_1, d_2) \Rightarrow d_1 = d_2$),
so Theorem~\ref{thm:bilemma} applies directly: no circuit explanation
can be simultaneously faithful and stable, regardless of whether
decisiveness is required.

\paragraph{Rashomon is derived.}
In the Lean formalization, the Rashomon property for circuits is
\emph{derived} from constructive finite types, not axiomatized.
We define two configurations (\texttt{circuitAlpha},
\texttt{circuitBeta}) computing the same function with different
decompositions; the Lean kernel verifies
$\mathrm{observe}(\texttt{circuitAlpha}) =
\mathrm{observe}(\texttt{circuitBeta})$ and
$\mathrm{incompatible}(\mathrm{explain}(\texttt{circuitAlpha}),
\mathrm{explain}(\texttt{circuitBeta}))$
by \texttt{decide} (zero axioms).

\paragraph{Empirical grounding.}
\citet{meloux2025circuits} found 85 distinct valid circuits with zero
circuit error for a simple XOR task, each admitting an average of 535.8
valid interpretations.
Less than 2\% of networks had a unique circuit mapping.
This extreme multiplicity confirms that the Rashomon property holds with
overwhelming force for circuit explanations in practice.

\paragraph{Resolution.}
The bilemma forecloses averaging-based resolutions because the circuit
explanation space has no neutral element.
The resolution requires \emph{enriching} the explanation space: report the
equivalence class of valid circuits---the computational structure shared
across \emph{all} valid decompositions---rather than any individual
decomposition.
This is the circuit analogue of the CPDAG in causal discovery: report what
is shared, acknowledge what is underdetermined.

\paragraph{Implication for AI safety.}
Safety cases that depend on identifying ``the circuit'' for a behavior are
subject to this impossibility: a circuit found in one training run may not
exist in a functionally equivalent run.

%% ----- Instance 2: Feature Attribution -----
\subsection{Instance 2: Feature Attribution (SHAP)}
\label{sec:attribution}

\paragraph{Setup.}
$\Theta$ is the space of trained models (e.g., GBDTs with different
random seeds);
$H$ is the space of feature importance rankings;
$Y$ is the space of predictions (loss-equivalent models map to the same
$y$).
Two rankings are incompatible if they order some feature pair in opposite
directions (feature~$j$ ranked above feature~$k$ vs.\ below).

\paragraph{The trilemma applies (not the bilemma).}
The ranking space admits \emph{ties} as a neutral element: a ranking that
assigns equal importance to features $j$ and $k$ is compatible with both
$j > k$ and $j < k$.
Because a neutral element exists, F+S is achievable (top-left cell of
Table~\ref{tab:tightness})---the bilemma does not apply.
However, ties sacrifice decisiveness, so the full trilemma still holds.

\paragraph{Resolution: DASH.}
DASH (Deterministic Averaged SHAP) averages SHAP values across an ensemble
of equivalent models, producing ties for feature pairs where the ensemble
disagrees.
These ties are the neutral elements that make F+S achievable.
DASH is provably optimal: by the Hunt--Stein theorem for the permutation
group $S_m$ acting on within-group features, the best equivariant estimator
is the Pitman estimator---the group-orbit average---which is exactly DASH.
Empirically, DASH resolves the instability on 68\% of 77 OpenML datasets
tested (under the standard workflow; the model-only component is ${\sim}63\%$
of the total variation).
The full quantitative treatment---$1/(1{-}\rho^2)$ divergence
rate, Spearman bounds, and model-specific
instantiations---appears in~\citet{anonymous2026attribution}.

%% ----- Instance 3: Causal Discovery -----
\subsection{Instance 3: Causal Discovery}
\label{sec:causal}

\paragraph{Setup.}
$\Theta$ is the space of DAGs in a Markov equivalence class;
$H$ is the space of fully oriented edge sets;
$Y$ is the conditional independence (CI) structure shared by all DAGs in
the class.
Two edge orientations are incompatible if they disagree on the direction of
at least one edge.

\paragraph{The Rashomon property.}
Markov equivalence is the canonical example of observational equivalence:
DAGs that encode the same CI structure are indistinguishable from
observational data.
For the minimal case (3-node chain $A \to B \to C$ vs.\ fork
$A \leftarrow B \to C$), same CI structure, different edge orientations.
In the Lean formalization, this is derived by \texttt{decide} (zero
axioms).

\paragraph{Impossibility.}
No stable orientation rule can be faithful to all datasets when the
equivalence class contains $\geq 2$ DAGs.
``Stable'' means the rule gives the same answer for all datasets consistent
with the CI structure; ``faithful'' means the answer matches the
data-optimal DAG.
Since different datasets can make different DAGs optimal, stability forces
a single answer, which cannot match all data-optimal DAGs simultaneously.

\paragraph{Resolution: CPDAG.}
The standard resolution in causal discovery is the CPDAG (completed
partially directed acyclic graph): orient edges that have the same
direction in all Markov-equivalent DAGs, leave the rest undirected.
This is precisely the neutral-element resolution: undirected edges are
compatible with both orientations.
The CPDAG achieves F+S by sacrificing decisiveness on the ambiguous edges.

%% ----- Instance 4: Attention Maps -----
\subsection{Instance 4: Attention Maps}
\label{sec:attention}

\paragraph{Setup.}
$\Theta$ is the space of transformer configurations (weight vectors after
different training runs);
$H$ is the space of token importance scores (or rankings) derived from
attention weights;
$Y$ is the space of input-output mappings.
\citet{damour2022underspecification} established that equivalent networks
(same loss on held-out data) can have markedly different internal
representations, including attention patterns.

\paragraph{The bilemma applies for argmax attention.}
When the explanation is the identity of the most-attended token
($H = \{\text{token}_1, \ldots, \text{token}_n\}$ with
$\mathrm{incompatible}(t_i, t_j) \Leftrightarrow i \neq j$),
the explanation space is maximally incompatible.
Theorem~\ref{thm:bilemma} applies: no method for identifying ``the most
important token'' can be simultaneously faithful and stable.

\paragraph{The trilemma applies for full attention distributions.}
When $H$ is the space of full attention distributions (probability vectors
over tokens), ties exist as neutral elements (uniform attention is
compatible with any native distribution under a suitable incompatibility
relation).
The trilemma applies but the bilemma does not; F+S is achievable via
averaging.

\paragraph{Empirical evidence.}
Using DistilBERT with 10 weight-perturbed variants ($\sigma = 0.02$),
88\% of adjacent token pairs exhibit flip rates exceeding 10\%.
Under actual fine-tuning (5 seeds, 1 epoch on SST-2), 14.5\% of pairs
are unstable---structural instability that persists across training
durations (identical at 1 and 3 epochs).
Mean Spearman correlation of token rankings across variants is 0.637,
comparable to single-model SHAP instability at moderate collinearity
($\rho \approx 0.5$).

% === PAGES 6-10 (merged from neurips_universal_part2.tex) ===
% neurips_universal_part2.tex — Sections 5–8 of NeurIPS 2026 paper
% "The Explanation Impossibility: A Universal Trilemma for Interpreting Underspecified Systems"
% This file is a LaTeX FRAGMENT. Insert after §4 in the main document.

%%%%%%%%%%%%%%%%%%%%%%%%%%%%%%%%%%%%%%%%%%%%%%%%%%%%%%%%%%%%%
%% §5. THE RESOLUTION
%%%%%%%%%%%%%%%%%%%%%%%%%%%%%%%%%%%%%%%%%%%%%%%%%%%%%%%%%%%%%

\section{The Resolution}
\label{sec:resolution}

The impossibility admits two qualitatively different resolutions, determined entirely by the structure of the explanation space~$H$.

\subsection{Case 1: Neutral Element Exists (Rankings, Continuous Scores)}

When $H$ contains a \emph{neutral element}---a value compatible with every native explanation (e.g., a tie in a ranking)---the $G$-invariant projection is uniquely optimal.

\begin{theorem}[Hunt--Stein optimality {\citep[Theorem~9.3.3]{lehmann2005testing}}]
\label{thm:hunt-stein}
For a finite group $G$ acting on the parameter space, the best equivariant estimator is the Pitman estimator: the average over the group orbit.
\end{theorem}

For feature attributions under the permutation group $S_m$ acting on $m$ within-group features, the Pitman estimator is the sample mean---precisely DASH consensus averaging.

\begin{definition}[DASH consensus attribution]
\label{def:dash}
Given $M$ independently trained models $f_1, \ldots, f_M$, the \emph{consensus attribution} for feature $j$ is $\bar{\varphi}_j = \frac{1}{M} \sum_{i=1}^{M} |\varphi_j(f_i)|$, where $\varphi_j(f_i)$ denotes the SHAP value of feature $j$ under model~$f_i$.
\end{definition}

DASH is a \emph{post-hoc explanation method}: the $M$ models are trained independently using the practitioner's existing pipeline; DASH only averages their SHAP values. No model is modified, no ensemble prediction is formed, and the deployed model remains unchanged.

\begin{corollary}[DASH equity]
\label{cor:equity}
For a balanced ensemble (each feature serves as first-mover in exactly $M/m$ models), $\bar{\varphi}_j = \bar{\varphi}_k$ for any two features $j, k$ in the same symmetry group.
\end{corollary}

\paragraph{Optimality.}
By the Cauchy--Schwarz inequality (Titu's lemma), any unbiased linear estimator satisfies $\Var(\hat{\mu}_j) \geq \sigma_j^2/M$, with equality iff $w_i = 1/M$ for all~$i$ (Lean-verified via \texttt{sum\_squares\_ge\_inv\_M}). DASH achieves this bound with $\Var(\bar{\varphi}_j) = \sigma_j^2/M$, making it the minimum-variance unbiased \emph{linear} estimator. By the Rao--Blackwell theorem (standard; not formalized in Lean), DASH is the minimum-variance unbiased estimator among \emph{all} estimators, since $\bar{\varphi}_j$ is a function of the sufficient statistic $\sum_i \varphi_j(f_i)$.

\paragraph{Ensemble size.}
For between-group stability $S_{jk} \geq 1 - \delta$ with importance gap~$\Delta_{jk}$ and attribution noise~$\sigma_{jk}$:
\begin{equation}
\label{eq:ensemble-size}
M_{\min} = \left\lceil \frac{\sigma_{jk}^2}{\Delta_{jk}^2} \cdot \bigl(\Phi^{-1}(1-\delta)\bigr)^2 \right\rceil.
\end{equation}
At $\delta = 0.05$: $M_{\min} = \lceil 2.71 \cdot \sigma_{jk}^2 / \Delta_{jk}^2 \rceil$. DASH achieves this bound with equality (to first order).

\subsection{Case 2: No Neutral Element (Circuits, SHAP Sign)}

For binary explanation spaces---$H = \{\text{positive}, \text{negative}\}$ (SHAP sign), $H = \{\text{selected}, \text{not selected}\}$ (feature selection), or circuit decompositions---no neutral element exists. Every answer is committal.

\begin{theorem}[No averaging resolution for binary $H$]
\label{thm:no-averaging}
When $H$ has no neutral element, faithful forces $E(\theta) = \operatorname{explain}(\theta)$ for all~$\theta$. Then $E$ is decisive, and by the trilemma, $E$ cannot be stable. No averaging-based fix (including DASH) applies.
\end{theorem}

The resolution requires \emph{enriching} the explanation space: replace $H$ with an enlarged space $\tilde{H} \supset H$ that contains equivalence classes as valid outputs. For circuit explanations, this means reporting the equivalence class of valid decompositions---the ``CPDAG of circuits''---rather than any individual decomposition. This parallels the CPDAG resolution in causal discovery: report what is shared across all valid structures, acknowledge what is underdetermined.

\paragraph{Practitioner diagnostic.}
The tightness classification is a design tool: \emph{check whether your explanation space has a neutral element}. If yes (rankings, continuous scores), apply DASH. If no (binary decisions, circuit decompositions), enrich $H$ or accept unfaithfulness at $\geq 50\%$ of Rashomon pairs.

\begin{center}
\small
\begin{tabular}{l|cc}
 & Neutral exists & No neutral \\
\hline
\textbf{Resolution} & DASH (averaging) & Enrich $H$ (equiv.\ classes) \\
\textbf{Cost} & $M \times$ training & Reduced decisiveness \\
\textbf{Guarantee} & Variance $\sigma^2/M$ & Stability + faithfulness \\
\end{tabular}
\end{center}


%%%%%%%%%%%%%%%%%%%%%%%%%%%%%%%%%%%%%%%%%%%%%%%%%%%%%%%%%%%%%
%% §6. EMPIRICAL VALIDATION
%%%%%%%%%%%%%%%%%%%%%%%%%%%%%%%%%%%%%%%%%%%%%%%%%%%%%%%%%%%%%

\section{Empirical Validation}
\label{sec:experiments}

We summarize four key results; the full experimental suite (18 experiments, 77 datasets) is in the supplement.

\subsection{Prevalence}

Of 77 public datasets (OpenML CC-18 + sklearn), 92\% have feature pairs with $|\rho| > 0.5$, and \textbf{52 (68\%)} exhibit attribution instability (flip rate $> 10\%$). The 95\% Wilson confidence interval is $[56\%, 77\%]$. For datasets with $P \geq 20$ features ($n = 43$), the rate rises to 93\% (40/43; Wilson CI: $[81\%, 98\%]$). With 20 models per dataset, the survey achieves only 32\% power at the 10\% threshold (60\% at 15\%), making 68\% a conservative lower bound; the true prevalence is likely 75--85\%.

\subsection{Cross-Implementation Consistency}

\begin{table}[t]
\centering
\caption{Attribution instability across three GBDT implementations on Breast Cancer (50 models each).}
\label{tab:cross-impl}
\small
\begin{tabular}{@{}lccc@{}}
\toprule
Metric & XGBoost & LightGBM & CatBoost \\
\midrule
Unstable pairs (flip $> 10\%$) & 162 & 183 & 160 \\
Max flip rate & 0.500 & 0.500 & 0.500 \\
Z-test $r(Z, \text{flip})$ & $-0.898$ & $-0.901$ & $-0.892$ \\
\bottomrule
\end{tabular}
\end{table}

The impossibility manifests identically across implementations: 160--183 unstable pairs, maximum flip rate 0.500 (the theoretical ceiling), and diagnostic correlation $|r| \geq 0.89$ in all three (Table~\ref{tab:cross-impl}). The instability is a property of sequential boosting under collinearity, not of any specific software.

\subsection{DASH Benchmark}

\begin{table}[t]
\centering
\caption{Flip rate comparison: DASH vs.\ alternatives. Synthetic Gaussian data ($P{=}10$, $m{=}5$, $N{=}1000$). 95\% bootstrap CIs.}
\label{tab:dash-benchmark}
\small
\begin{tabular}{@{}lcccc@{}}
\toprule
Method & $\rho = 0.5$ & $\rho = 0.7$ & $\rho = 0.9$ & Pairs resolved \\
\midrule
Single model       & 16.0\% [13.8, 18.1] & 16.4\% [14.3, 18.5] & 18.7\% [16.4, 21.1] & 40/40 \\
Bootstrap SHAP     & 16.2\% [13.8, 18.2] & 16.1\% [14.2, 18.0] & 18.8\% [16.4, 21.4] & 40/40 \\
Subsampled SHAP    & 16.3\% [13.8, 18.5] & 16.1\% [14.2, 17.9] & 18.8\% [16.5, 21.3] & 40/40 \\
\textbf{DASH ($M{=}25$)} & \textbf{3.8\%} [2.8, 4.9] & \textbf{3.1\%} [2.4, 3.9] & \textbf{3.7\%} [2.7, 4.7] & 40/40 \\
CI SHAP            & 0.2\% [0.1, 0.4]    & 0.8\% [0.6, 1.1]    & 3.1\% [2.1, 4.6]    & $\sim$22/40 \\
\bottomrule
\end{tabular}
\end{table}

DASH dominates all alternatives at every $\rho$ when measured on resolved pairs (Table~\ref{tab:dash-benchmark}). Bootstrap and subsampled SHAP address estimation noise (a single-model phenomenon), not the Rashomon effect (a multi-model phenomenon): their flip rates match the single-model baseline. CI SHAP achieves lower rates only by refusing to rank ${\sim}45\%$ of pairs---the impossibility theorem in action. The two approaches correspond to the two families in the Design Space Theorem: DASH (Family $\mathcal{B}$: stable, faithful, with ties) and CI SHAP (Family $\mathcal{A}$: explicitly incomplete).

\subsection{Monte Carlo Validation}

The theoretical flip rate formula $\Phi(-\text{SNR})$ matches empirical rates within 1.35 percentage points across 10,000 trials at each of 8 correlation levels ($\rho \in \{0.1, 0.3, 0.5, 0.7, 0.8, 0.9, 0.95, 0.99\}$). The maximum discrepancy occurs at $\rho = 0.5$ (empirical 39.9\% vs.\ theoretical 38.6\%). The SNR calibration on 1,325 feature pairs from 6 datasets achieves $R^2 = 0.94$; all 228 pairs with $\text{SNR} > 1.96$ have flip rate $< 5\%$.


%%%%%%%%%%%%%%%%%%%%%%%%%%%%%%%%%%%%%%%%%%%%%%%%%%%%%%%%%%%%%
%% §7. RELATED WORK
%%%%%%%%%%%%%%%%%%%%%%%%%%%%%%%%%%%%%%%%%%%%%%%%%%%%%%%%%%%%%

\section{Related Work}
\label{sec:related}

\paragraph{Attribution impossibilities.}
\citet{bilodeau2024impossibility} prove that completeness and linearity cannot coexist for feature attributions---a constraint on the \emph{method} axis. Our result operates on a different axis: cross-model \emph{stability} under collinearity. The Bilodeau impossibility constrains which axioms a method can satisfy for a fixed model; ours constrains what any method can achieve across the Rashomon set. The two results are complementary: Bilodeau's applies even without collinearity; ours applies even to methods that are not linear. Our result additionally provides quantitative architecture-discriminating bounds ($1/(1-\rho^2)$ for GBDT, $\infty$ for Lasso), a constructive resolution (DASH), and machine verification (335 Lean~4 theorems).

\paragraph{Fairness impossibilities.}
\citet{chouldechova2017fair} and \citet{kleinberg2017inherent} prove that calibration, balance, and equal false positive rates cannot coexist when base rates differ. The structural parallel is precise: three desirable properties conflict, and the resolution requires sacrificing one. Our trilemma (faithfulness, stability, completeness) extends this pattern to the explanation domain, and our tightness classification provides a diagnostic for which sacrifice is necessary.

\paragraph{Underspecification and model multiplicity.}
\citet{damour2022underspecification} demonstrated that ML pipelines are underspecified: models with equivalent held-out performance behave differently on deployment-relevant criteria. Our impossibility is the explanation-level consequence of this phenomenon. Where D'Amour et~al.\ document that predictions vary, we prove that feature \emph{rankings} must vary---and characterize exactly when, by how much, and what the complete design space looks like. The Rashomon literature---\citet{breiman2001random}, \citet{fisher2019models}, \citet{semenova2022existence}---established model multiplicity as an empirical and theoretical concern; our theorem is the explanation-theoretic capstone, proving that the Rashomon property is \emph{sufficient} for attribution impossibility (zero additional axioms).

\paragraph{Circuit multiplicity.}
\citet{meloux2025circuits} found 85 distinct valid circuits with zero circuit error for a simple XOR task, each admitting an average of 535.8 valid interpretations. Our impossibility is the formal consequence: the bilemma (Theorem~\ref{thm:bilemma}) proves that no circuit explanation can be simultaneously faithful and stable, and the tightness classification shows that the resolution requires enriching the explanation space (equivalence classes) rather than averaging.

\paragraph{Formalization methodology.}
\citet{nipkow2009social} formalized Arrow's impossibility theorem in Isabelle/HOL. Our formalization is the first machine-verified impossibility result in explainable AI. Both demonstrate that formalization catches errors missed by informal proof (we found two logical inconsistencies). The Arrow and attribution impossibilities share surface structure (three desirable properties conflict) but differ fundamentally: Arrow aggregates many inputs; the attribution impossibility maps one input to one output.

\paragraph{The disagreement problem.}
\citet{krishna2022disagreement} document that different explanation methods produce different outputs for the same model---an empirical observation. Our result proves the complementary phenomenon: the \emph{same} method produces different outputs for different equivalent models, and this is a mathematical necessity, not an empirical accident.

\begin{table}[t]
\centering
\caption{Positioning relative to prior work. Comparison criteria reflect the specific contributions of the present work.}
\label{tab:prior-work}
\small
\begin{tabular}{@{}lcccc@{}}
\toprule
Result & Stability? & Quantitative? & Resolution? & Formal? \\
\midrule
Bilodeau et al.\ (2024) & No & No & No & No \\
D'Amour et al.\ (2022) & Empirical & No & No & No \\
Meloux et al.\ (2025) & Empirical & No & No & No \\
Krishna et al.\ (2022) & Empirical & No & No & No \\
Chouldechova (2017) & (Fairness) & Yes & No & No \\
\textbf{This paper} & \textbf{Impossibility} & \textbf{Yes} & \textbf{DASH} & \textbf{Lean 4} \\
\bottomrule
\end{tabular}
\end{table}


%%%%%%%%%%%%%%%%%%%%%%%%%%%%%%%%%%%%%%%%%%%%%%%%%%%%%%%%%%%%%
%% §8. DISCUSSION
%%%%%%%%%%%%%%%%%%%%%%%%%%%%%%%%%%%%%%%%%%%%%%%%%%%%%%%%%%%%%

\section{Discussion}
\label{sec:discussion}

The central insight of this paper is that \emph{the ceiling of explainability depends on the explanation space}.

\paragraph{Binary spaces are fundamentally harder.}
For binary explanation spaces (circuit decompositions, SHAP sign, feature selection), no averaging fix exists. The bilemma proves that faithful + stable alone is impossible---even dropping completeness does not help. The only resolution is to \emph{enrich} $H$ with equivalence classes, trading decisiveness for stability and faithfulness. This has direct consequences for AI safety: circuit-based safety cases that depend on identifying ``the circuit'' for a behavior are subject to this impossibility. Safety claims built on single-network circuit analysis are as unstable as single-model SHAP rankings. The remedy is multi-network verification: verify circuit-level claims across multiple equivalent networks, or qualify them as instance-specific.

\paragraph{Graded spaces admit averaging.}
For ranking and continuous-score explanation spaces, DASH consensus averaging achieves the Cauchy--Schwarz variance bound with Hunt--Stein optimality. The ensemble size formula $M_{\min} = \lceil 2.71 \sigma^2/\Delta^2 \rceil$ gives practitioners a concrete budget. The tightness table (neutral element $\times$ committal element) serves as a design tool: before building an XAI system, check which cell your explanation space occupies to know which resolution applies.

\paragraph{Regulatory implications.}
Single-model feature explanations are unreliable under collinearity, affecting regulatory compliance in settings where explanation stability is required (e.g., EU AI Act Art.~13, US SR~11-7). Ensemble explanations provide a defensible alternative.

\paragraph{Limitations.}
Several assumptions warrant discussion.
(i)~The equicorrelation structure (uniform $\rho$ within groups) simplifies the axioms; the Rashomon property holds pairwise for arbitrary correlation structures.
(ii)~The mechanistic interpretability instance rests on \citet{meloux2025circuits}, which studies a single task (XOR with 3-bit inputs); extending to larger architectures and tasks is important future work.
(iii)~Circuit equivalence classes are not yet formalized in Lean; the bilemma and unfaithfulness bound are verified, but the enrichment resolution is argued informally.
(iv)~The global proportionality axiom has empirical CV of 0.35--0.66; under variable proportionality constants, DASH achieves approximate rather than exact equity, with the equity violation bounded by the CV of $c$ across models.
(v)~The prevalence survey achieves only 32\% power at the 10\% flip rate threshold; the reported 68\% is a conservative lower bound.

\paragraph{Lean formalization.}
The impossibility, resolution, and design space are machine-verified in Lean~4: 335 theorems, 0~\texttt{sorry}, 58 files. The core impossibility (Theorem~\ref{thm:trilemma}) has \emph{zero} behavioral axiom dependencies---only the Rashomon property as a hypothesis. The mechanistic interpretability instance derives the Rashomon property from constructive finite types (zero axioms via \texttt{decide}). Beyond certifying correctness, the formalization served as a methodology: it caught two logical inconsistencies and one type mismatch that survived informal review. We recommend formalizing impossibility theorems as a general practice in ML theory.

\paragraph{Broader impact.}
Feature attribution methods are used to make consequential decisions: which clinical interventions to pursue, which loan applicants to approve, which model behaviors to trust. Our theorem establishes that single-model explanations of collinear features are unreliable---a ``known and foreseeable circumstance'' that practitioners, deployers, and regulators should account for. The constructive resolution (DASH for graded spaces, equivalence classes for binary spaces) provides actionable remedies. We encourage practitioners to adopt ensemble explanations and to report instability diagnostics alongside feature rankings, so that the consumers of explanations---patients, applicants, auditors---receive honest assessments of explanation reliability.


\bibliographystyle{plainnat}
\bibliography{references}

\end{document}
